
\exo

Sans calcul, comparer les nombres suivants.

\vspace{-6pt}
\setlength{\columnsep}{1cm}
\setlength{\columnseprule}{0pt}
\begin{multicols}{3}
\begin{enumerate}
\item $\dfrac{1}{2}$\; et \; $\dfrac{1}{3}$.
\medskip
\item $\dfrac{1}{-3}$\; et \; $\dfrac{1}{-2}$.
\item $\dfrac{1}{2,05}$\; et \; $\dfrac{1}{1,95}$.
\medskip
\item $\dfrac{1}{3}$\; et \; $0,5$.
\item $\dfrac{1}{5+\sqrt{2}}$\; et \; $\dfrac{1}{5-\sqrt{2}}$.
\medskip
\item $\dfrac{1}{-8}$\; et \; $\dfrac{1}{3}$.
\end{enumerate}
\end{multicols}

\exo

Soit $x$ un nombre réel. Donner un encadrement de $\frac{1}{x}$ dans chacun des cas suivants.
\vspace{-6pt}
\setlength{\columnsep}{1cm}
\setlength{\columnseprule}{0pt}
\begin{multicols}{3}
\begin{enumerate}
\item $2 \leqslant x \leqslant 5$.
\medskip
\item $-7 \leqslant x \leqslant -1$.
\item $0 < x < 3$.
\medskip
\item $-10 \leqslant x < 0$.
\item $x\in\intfo{-10}{0}\cup\intof{0}{10}$.
\medskip
\item $x\in\intfo{-10}{0}\cup\intoo{0}{+\infty}$.
\end{enumerate}
\end{multicols}

\exo

\begin{minipage}[t]{11.5cm}
On a tracé ci-contre la courbe $\cC$ représentative de la fonction inverse sur l'intervalle $\intof{0}{4}$. La droite $d$ est constituée des points d'ordonnée $2$.
\begin{enumerate}
\item Résoudre graphiquement l'équation $\frac{1}{x}=2$

 sur l'intervalle $\intof{0}{4}$.
\item
\begin{enumerate}
\item Mettre en évidence sur le graphique les points de $\cC$ qui sont en-dessous de la droite $d$.
\item En déduire les solutions dans $\intof{0}{4}$ de l'inéquation $\frac{1}{x}<2$.
\end{enumerate}
\end{enumerate}
\end{minipage}
\hfill
\begin{minipage}[t]{5.5cm}
~

\vspace{-12pt}
\def\xmin{-1}
\def\xmax{5}
\def\ymin{-1}
\def\ymax{5}
\def\xunite{0.8cm}
\def\yunite{0.8cm}
\psset{unit=\xunite, xunit=\xunite, yunit=\yunite, algebraic=true, dimen=middle, dotstyle=*, dotsize=3pt 0, linewidth=0.8pt, arrowsize=3pt 2, arrowinset=0.25, comma=true}
\begin{pspicture*}(\xmin,\ymin)(\xmax,\ymax)
\psgrid[xunit=\xunite, yunit=\yunite, subgriddiv=0, gridlabels=0, gridcolor=lightgray](0,0)(\xmin,\ymin)(\xmax,\ymax)
\psaxes[labelFontSize=\scriptstyle, xAxis=true, yAxis=true, Dx=1, Dy=1, ticksize=-2pt 0, subticks=1]{->}(0,0)(-0.99,-0.99)(\xmax,\ymax)[{\footnotesize $x$},135][{\footnotesize $y$},-45]
\psplot[linewidth=1.2pt, plotpoints=200]{0.1}{4}{1/x}
\psplot[linestyle=dashed, dash=2pt 2pt,linewidth=1.2pt, plotpoints=200]{\xmin}{\xmax}{2}
\begin{footnotesize}
\uput{4pt}[045](0.25,4){$\cC$}
\uput{4pt}[090](3.5,2){$d$}
\end{footnotesize}
\end{pspicture*}
\end{minipage}

\exo

\begin{minipage}[t]{7.5cm}
En s'aidant de la représentation graphique de la fonction inverse rappelée ci-contre, résoudre dans $\R^*$ les équations et inéquations proposées :

\vspace{-6pt}
\setlength{\columnsep}{1cm}
\setlength{\columnseprule}{0pt}
\begin{multicols}{2}
\begin{enumerate}
\item $\frac{1}{x}=4$.
\medskip
\item $\frac{1}{x}<4$.
\medskip
\item $\frac{1}{x}\geqslant 4$.
\medskip
\item $\frac{1}{x}=0,5$.
\medskip
\item $\frac{1}{x}\leqslant 0,5$.
\medskip
\item $\frac{1}{x}>0,5$.
\item $\frac{1}{x}=-3$.
\medskip
\item $\frac{1}{x}<-3$.
\medskip
\item $\frac{1}{x}\geqslant -3$.
\medskip
\item $\frac{1}{x}\leqslant -0,1$.
\medskip
\item $\frac{1}{x}\geqslant -0,1$.
\medskip
\item $\frac{1}{x}>-0,1$.
\end{enumerate}
\end{multicols}
\vspace{-6pt}
\end{minipage}
\hfill
\begin{minipage}[t]{8cm}
~

\vspace{-12pt}
\def\xmin{-4}
\def\xmax{4}
\def\ymin{-4}
\def\ymax{4}
\def\xunite{1.0cm}
\def\yunite{1.0cm}
\psset{unit=\xunite, xunit=\xunite, yunit=\yunite, algebraic=true, dimen=middle, dotstyle=*, dotsize=3pt 0, linewidth=0.8pt, arrowsize=3pt 2, arrowinset=0.25, comma=true}
\begin{pspicture*}(\xmin,\ymin)(\xmax,\ymax)
%\psgrid[xunit=\xunite, yunit=\yunite, subgriddiv=0, gridlabels=0, gridcolor=lightgray](0,0)(\xmin,\ymin)(\xmax,\ymax)
\psaxes[labels=none,labelFontSize=\scriptstyle, xAxis=true, yAxis=true, Dx=10, Dy=10, ticksize=-2pt 0, subticks=1]{->}(0,0)(\xmin,\ymin)(\xmax,\ymax)[{\footnotesize $x$},135][{\footnotesize $y$},-45]
\psplot[linewidth=1.2pt, plotpoints=200]{\xmin}{-0.1}{1/x}
\psplot[linewidth=1.2pt, plotpoints=200]{0.1}{\xmax}{1/x}
\end{pspicture*}
\end{minipage}

\exo

On note $f$ la fonction inverse, $a$ et $b$ deux réels non nuls.

Les affirmations suivantes sont-elles vraies ou fausses ? Justifier la réponse.

\begin{enumerate}
\item On a toujours $f(f(a))=a$.
\smallskip
\item Si $f(a)=b$, alors $f(b)=a$.
\smallskip
\item Si $f(a)=a$, alors $a=1$.
\end{enumerate}

\exo

Résoudre les équations suivantes :

\vspace{-6pt}
\setlength{\columnsep}{1cm}
\setlength{\columnseprule}{0pt}
\begin{multicols}{3}
\begin{enumerate}
\item $\dfrac{1}{x^2}=4$.
\medskip
\item $\dfrac{1}{x^3}=4$.
\medskip
\item $\dfrac{1}{\sqrt{x}}=4$.
\medskip
\item $\dfrac{3}{x^2}=9$.
\medskip
\item $\dfrac{-3}{x^3}=9$.
\medskip
\item $\dfrac{20}{\sqrt{x}}=5$.
\medskip
\item $\dfrac{4}{x-1}=2$.
\medskip
\item $-\dfrac{7}{3x+6}=4$.
\medskip
\item $\dfrac{-2}{-3x+9}=8$.
\medskip
\end{enumerate}
\end{multicols}
\vspace{-6pt}

\exo

Construire le tableau de signe des expressions suivantes :

\vspace{-6pt}
\setlength{\columnsep}{1cm}
\setlength{\columnseprule}{0pt}
\begin{multicols}{3}
\begin{enumerate}
\item $Q(x)=\dfrac{x+7}{5x-10}$.
\medskip
\item $Q(x)=\dfrac{-2x+6}{x+8}$.
\medskip
\item $Q(x)=\dfrac{6-3x}{4x+12}$.
\medskip
\item $Q(x)=\dfrac{(x+7)(x-5)}{5x-10}$.
\medskip
\item $Q(x)=\dfrac{-x+3}{(x-2)(5-3x)}$.
\medskip
\item $Q(x)=\dfrac{x^3}{3x-12}$.
\medskip
\end{enumerate}
\end{multicols}
\vspace{-6pt}

\exo

Résoudre dans $\R$ les inéquations suivantes :

\vspace{-6pt}
\setlength{\columnsep}{1cm}
\setlength{\columnseprule}{0pt}
\begin{multicols}{3}
\begin{enumerate}
\item $\dfrac{x+7}{5x-10}>0$.
\medskip
\item $\dfrac{-2x+6}{x+8}<0$.
\medskip
\item $\dfrac{6-3x}{4x+12}\geqslant 0$.
\medskip
\item $\dfrac{(x+7)(x-5)}{5x-10}\geqslant 0$.
\medskip
\item $\dfrac{-x+3}{(x-2)(5-3x)}\leqslant 0$.
\medskip
\item $\dfrac{x^3}{3x-12}<0$.
\medskip
\end{enumerate}
\end{multicols}
\vspace{-6pt}

\exo

\begin{minipage}[t]{11.5cm}
On a tracé ci-contre la courbe $\cC$ représentative de la fonction inverse, ainsi que la droite $\cD$ d'équation $y=x$.
\begin{enumerate}
\item En observant ce dessin, Alice affirme :
\[x\geqslant \tfrac{1}{x}\qeqq x\in\intfo{-1}{0}\cup\intfo{1}{+\infty}.\]
A-t-elle raison ? Justifier.
\item Montrer que, pour tout réel $x$ non nul :
\[x-\tfrac{1}{x}=\tfrac{x^2-1}{x}.\]
\item Prouver la conjecture d'Alice en étudiant le signe de $\frac{x^2-1}{x}$.
\end{enumerate}
\end{minipage}
\hfill
\begin{minipage}[t]{5.5cm}
~

\vspace{-12pt}
\def\xmin{-3}
\def\xmax{3}
\def\ymin{-3}
\def\ymax{3}
\def\xunite{0.8cm}
\def\yunite{0.8cm}
\psset{unit=\xunite, xunit=\xunite, yunit=\yunite, algebraic=true, dimen=middle, dotstyle=*, dotsize=3pt 0, linewidth=0.8pt, arrowsize=3pt 2, arrowinset=0.25, comma=true}
\begin{pspicture*}(\xmin,\ymin)(\xmax,\ymax)
\psgrid[xunit=\xunite, yunit=\yunite, subgriddiv=0, gridlabels=0, gridcolor=lightgray](0,0)(\xmin,\ymin)(\xmax,\ymax)
\psaxes[labelFontSize=\scriptstyle, xAxis=true, yAxis=true, Dx=1, Dy=1, ticksize=-2pt 0, subticks=1]{->}(0,0)(-2.99,-2.99)(\xmax,\ymax)[{\footnotesize $x$},135][{\footnotesize $y$},-45]
\psplot[linewidth=1.2pt, plotpoints=200]{\xmin}{-0.1}{1/x}
\psplot[linewidth=1.2pt, plotpoints=200]{0.1}{\xmax}{1/x}
\psplot[linestyle=dashed, dash=2pt 2pt,linewidth=1.2pt, plotpoints=200]{\xmin}{\xmax}{x}
\begin{footnotesize}
\uput{4pt}[090](2,0.5){$\cC$}
\uput{4pt}[135](2,2){$\cD$}
\end{footnotesize}
\end{pspicture*}
\end{minipage}

\exo

Alice loue une camionnette pour son déménagement. Le prix de la location est $75$ \euro{}, auquel s'ajoute le carburant à $1,50$ \euro{} le litre. La camionnette consomme $7$ litres de carburant aux $100$ kilomètres.
\begin{enumerate}
\item
\begin{enumerate}
\item Quel est le coût pour un trajet de $80$ kilomètres ?
\item Quel est le coût moyen de chaque kilomètre parcouru lors de ce trajet ?
\end{enumerate}
\item Plus généralement, on note $x$ le nombre de kilomètres parcourus.
\begin{enumerate}
\item Exprimer en fonction de $x$ le coût total $C(x)$ du trajet, en euros.
\item Exprimer en fonction de $x$ le coût moyen $C_M(x)$ d'un kilomètre, pour un trajet de $x$ km.
\item Combien de kilomètres faut-il parcourir pour que le coût moyen par kilomètre devienne inférieur à $2,40$ \euro{} ?
\end{enumerate}
\end{enumerate}

